\ifHAFOPedagogy
  \section{Escalas y convenciones}
  \label{sec:master-scope}
\else
  \section{Fundamentos de la comparación y sistemas estudiados}
  \label{sec:master-scope}

  La localización combina realizaciones de Lur'e, balance armónico,
  continuación e integración causal. La clasificación exige estudiar la
  atracción y las cuencas en el espacio de fases correspondiente. La geometría
  proporciona semillas y regiones de búsqueda; la topología caracteriza
  conjuntos invariantes, fronteras y condiciones de existencia.

  Las afirmaciones se distinguen por su fundamento algebraico, dinámico o
  topológico:
\fi

\begin{longtable}{@{}L{0.16\linewidth}L{0.28\linewidth}L{0.48\linewidth}@{}}
\caption{Alcance de los resultados algebraicos y numéricos.}\label{tab:master-levels}\\
\toprule
Nivel & Nombre & Interpretación\\
\midrule
\endfirsthead
\toprule
Nivel & Nombre & Interpretación\\
\midrule
\endhead
E0 & referencia bibliográfica & Se dispone de ecuaciones, parámetros o
resultados publicados.\\
E1 & desarrollo algebraico & Se han calculado identidades, determinantes,
transferencias, equilibrios, jacobianos, raíces o residuos.\\
E2 & semilla & Existe una condición inicial producida por función descriptiva,
familia auxiliar de Machado, punto perpetuo, mapa de conmutación, órbita de
ganancia cero o búsqueda acotada declarada.\\
E3 & continuación & La semilla fue transportada hasta el campo objetivo mediante
un protocolo de continuación reproducible.\\
E4 & dinámica & El conjunto final fue integrado y clasificado mediante
retornos, Lyapunov, prueba \(0\)--\(1\), Poincaré u otros diagnósticos finitos.\\
E5 & separación local & Condiciones iniciales predeclaradas alrededor de todos
los equilibrios no alcanzaron el candidato, o se registró el primer contacto.\\
E6 & reproducción independiente & Otro código, lenguaje o integrador repitió
la localización o la clasificación sin leer la trayectoria objetivo.\\
\bottomrule
\end{longtable}

\ifHAFOdigitalcontent
E5 expresa evidencia finita y no una prueba global. Si se ensayan \(N\)
condiciones sin observar contacto, el resultado exacto es «cero contactos bajo
ese muestreo, radios, horizonte y tolerancia». Las divergencias se registran
por separado y no cuentan como confirmación de separación.

La escala E0--E6 describe el máximo alcanzado por la evaluación integral de un
caso. La escala EV--TG0--EV--TG7 usada en la campaña
geométrico--topológica describe, en cambio, una ruta o una semilla concreta:
EV--TG2 corresponde a destino dinámico finito, EV--TG3 añade una obstrucción
local de equilibrio y los niveles superiores exigen evidencia de borde,
estructura exterior o certificación topológica. Las dos escalas no se
identifican término a término y no deben combinarse tomando el máximo de una
ruta como estado de todo el caso.
\fi

\subsection{Sistemas estudiados}

\begin{longtable}{@{}L{0.22\linewidth}C{0.08\linewidth}L{0.19\linewidth}L{0.43\linewidth}@{}}
\caption{Casos algebraicos, dinámicos y candidatos.}\label{tab:master-cases}\\
\toprule
Sistema o ruta & Orden & Máximo nivel & Estado según la evidencia\\
\midrule
\endfirsthead
\toprule
Sistema o ruta & Orden & Máximo nivel & Estado según la evidencia\\
\midrule
\endhead
Chua PWL canónico & \(q=1\) & E6 & Candidato caótico compatible con ocultedad
bajo dos campañas independientes de 504 sondas, ambas con 0 contactos. Sus
particiones de destinos no son comparables porque usan coordenadas, horizontes
y clasificadores distintos. Semilla DF verificada.\\
Chua PWL «nearby» & \(q=0.9998\) & E5 & Control rechazado: 1305
contactos de 4050 y 1395 divergencias; autoexcitado bajo el protocolo ensayado.\\
Chua PWL exacto de Danca, centrado & \(q=0.9998\) & E2 & Dos semillas
armónicas; no confundir con el extremo sesgado.\\
Chua PWL exacto de Danca, sesgado & \(q=0.9998\) & E5 & Extremo
regular/periódico, no caótico; 0/7200 contactos para \(r\leq 0.01\) y primer
contacto macroscópico a \(r=0.3\).\\
Chua arctangente de Wu & \(q=0.99\) & E4 & Reproducción bibliográfica ADM;
señal regular y sin equivalencia demostrada con un PVI Caputo causal de memoria
completa.\\
Chua arctangente c590 & \(q=0.9999\) & E5 & Candidato aperiódico bajo
diagnósticos finitos; \(0/8\,400\) contactos hasta \(r=0.3\),
\(22/2\,400\) en \(r=1\) y \(588/2\,700\) en \(r=2\), con 131
divergencias en este último bloque y 610 contactos totales. Su semilla procede
de búsqueda entera acotada y
refinamiento Caputo, no de función descriptiva.\\
Kalman--Fitts & \(q=1\) & E6 & La ruta DF directa está algebraicamente
obstruida; mapa de conmutación y homotopía signo--tanh producen un ciclo
estable compatible con ocultedad bajo \(0/48\) contactos finitos.\\
MAVPD, \(\xi=3.1\) & \(q=1\) & E6 & Ciclo periódico externo, 0/108
contactos. La etiqueta publicada de cuasiperiodicidad no se reproduce. Dos PP
exactos llegan a las órbitas simétricas en el experimento de puntos perpetuos.\\
MAVPD, \(\xi=3.5\) & \(q=1\) & E6 & Control negativo: aparecen 24
contactos; no es oculto bajo el protocolo común.\\
MAVPD local, \(\xi=2.85\) & \(q=1\) & E6 & Candidato compatible con caos bajo
diagnósticos finitos en \(\gamma=0.15380379839949113\); 0/112 contactos y
verificación Python--Julia.\\
PLL lead--lag & \(q=1\) & E6 & Ruta directa obstruida; órbita corredora en
ganancia cero y continuación a \(L=500\); 0/96 contactos cilíndricos.\\
Lorenz generalizado de Danca & \(q=0.995\) & E2 & Realización MIMO,
equilibrios, espectros y dos semillas FPP--A; existencia global,
absorción y precompacidad de perfiles demostradas. La dinámica desde las
semillas permanece sin caracterizar.\\
Rabinovich--Fabrikant de Danca & \(q=0.998\) & E2 & Realización MIMO,
cinco equilibrios, estabilidad de Matignon y ocho semillas FPP--A;
la localización numérica del atractor permanece pendiente.\\
Muñoz--Pacheco, piloto PP & \(q=0.97\) & E2 & Cuatro semillas FPP--A exactas y
trayectorias PECE de historia completa finitas hasta \(t=35\); sin
clasificación de atractor ni evento FPP--H convergente.\\
Candidatos adicionales & mixto & E0--E1 & Diseños o candidatos compatibles
con pruebas de estructura y un plan de validación L0--L6; su estado no corresponde a
atractores reproducidos.\\
\bottomrule
\end{longtable}

\ifHAFOdigitalcontent
Las 18 semillas y las tres ramas supervivientes del Chua fraccionario
pertenecen al mismo modelo; su número no determina la cantidad de atractores.
Las pruebas de DK2018, Fischer y EFORK3 evalúan exponentes o integradores; su alcance
no incluye evidencia de localización u ocultedad.
El nivel de la tabla es el máximo alcanzado por la evaluación integral. Cada
semilla, rama o trayectoria conserva su nivel propio y no hereda
automáticamente el máximo del caso.
\fi

\subsection{Convención de realimentación unificada}

Se adopta la convención
\begin{equation}
 {}^{\mathrm C}D_{t_0}^{q}x=Px+b\,\psi(r^{\mathsf T}x),
 \qquad
 W_+(z)=r^{\mathsf T}(zI-P)^{-1}b,
 \qquad z=s^q,
 \label{eq:master-wplus}
\end{equation}
y el cierre
\begin{equation}
 1-W_+(z)N=0.
 \label{eq:master-closure}
\end{equation}
La convención alternativa es
\begin{equation}
 W_-(z)=r^{\mathsf T}(P-zI)^{-1}b=-W_+(z),
 \qquad 1+W_-(z)N=0.
 \label{eq:master-wminus}
\end{equation}
Por tanto, \(N=1/\Re W_+\) y \(N=-1/\Re W_-\) son la misma condición. Esta
igualdad establece la equivalencia entre las dos convenciones de signo. Se
reserva \(q\) para el orden fraccionario y \(b\) para el vector de entrada.

\subsection{Interpretación de las magnitudes}

La amplitud de la semilla canónica de Chua entero es
\(A_0=5.856145086257\). Las transferencias del sistema original y de su
linealización alrededor de un sesgo se distinguen mediante el punto de
referencia; la pendiente local interviene una sola vez en la transferencia
incremental. En la construcción de Machado, los residuos del cierre auxiliar
con \(N^\mu\) y del balance físico con \(N\) se evalúan por separado.

El espectro de amplitud expresa la magnitud de la transformada de Fourier.
La densidad espectral de potencia expresa potencia por unidad de frecuencia
y exige la normalización correspondiente. La clasificación regular del
extremo PWL sesgado y el carácter aperiódico del candidato c590 se interpretan
junto con sus retornos, espectros y controles de integración.

La continuación Caputo conserva una historia de Volterra. La ecuación
\(J_f(p)f(p)=0\) define una semilla de arranque Picard--Caputo; el evento
secuencial FPP--H se evalúa sobre la trayectoria con su historia. La
compatibilidad numérica con ocultedad depende de los radios, el muestreo,
el horizonte y las tolerancias. Una demostración de ocultedad requiere
establecer el atractor y la separación de su cuenca respecto de vecindades
de los equilibrios. Si no existen equilibrios, esta última condición es
vacía, aunque la existencia y atracción siguen siendo necesarias.

\subsection{Fundamentos bibliográficos}

La forma de Lur'e, el balance armónico, la función descriptiva de Machado
y los puntos perpetuos enteros proporcionan métodos complementarios de
generación de semillas. Las formulaciones FPP--A y FPP--H de la
\cref{sec:perpetual-points} distinguen el arranque de Picard--Caputo de un
evento dependiente de la historia. Su utilidad se evalúa mediante integración
causal, comparación entre métodos y estudio de cuencas.
